\chapter{Entropy and Asymptotic Bounds}
%%%%%%%%%%%%%%%%%%%%%%%%%%%%%%%%%%%%%%%%%%%%%%%%%%%%%%%%%%%%%%%%%%%%%%%%%%%%%%%
\section{Overview}

This chapter develops the asymptotic relationship between Hamming ball sizes and the
$q$-ary entropy function $H_q$, and provides the key binomial lower bound used in the
Gilbert--Varshamov argument.

%%%%%%%%%%%%%%%%%%%%%%%%%%%%%%%%%%%%%%%%%%%%%%%%%%%%%%%%%%%%%%%%%%%%%%%%%%%%%%%
\section{Asymptotic upper bound on ball size}

\begin{theorem}[Entropy upper bound on Hamming balls]
\lean{ErrorCorrectingCodes.Codeword.hamming_ball_size_asymptotic_upper_bound}
\label{ErrorCorrectingCodes.Codeword.hamming_ball_size_asymptotic_upper_bound}
\leanok
\proofsource{ada-coding-book}{pdf-7251cb8e4bc5-p071-b001}
\uses{ErrorCorrectingCodes.Codeword.hamming_ball_size, ErrorCorrectingCodes.Codeword.qaryEntropy}
\difficulty{5}
Let $\alpha$ be a finite alphabet with at least two symbols, write $q = \abs{\alpha}$
for its size, and let $n$ be a natural number. Fix a real number $p$ with $0 < p \le 1 -
1/q$, and let
\[
H_q(p) \;=\; p\log_q(q-1) - p\log_q p - (1-p)\log_q(1-p)
\]
be the $q$-ary entropy function. Then for every codeword $c \in \alpha^n$, the Hamming
ball of radius $\lfloor np\rfloor$ about $c$, namely $B_{\lfloor np\rfloor}(c) = \{c'
\in \alpha^n : d(c',c) \le \lfloor np\rfloor\}$, satisfies
\[
\abs{B_{\lfloor np\rfloor}(c)} \;\le\; q^{\,H_q(p)\,n}.
\]
\end{theorem}

%%%%%%%%%%%%%%%%%%%%%%%%%%%%%%%%%%%%%%%%%%%%%%%%%%%%%%%%%%%%%%%%%%%%%%%%%%%%%%%
\section{Entropy algebra lemmas}

\begin{lemma}[Closed form for $q^{H_q(p)}$]
\lean{ErrorCorrectingCodes.Codeword.q_pow_qary_entropy_simp}
\label{ErrorCorrectingCodes.Codeword.q_pow_qary_entropy_simp}
\leanok
\uses{ErrorCorrectingCodes.Codeword.qaryEntropy}
\difficulty{4}
Let $q \ge 2$ be an integer and let $p$ be a real number with $0 < p < 1$. Then the
$q$-th power of the $q$-ary entropy $H_q(p)$ is given by
\[
q^{H_q(p)} \;=\; (q-1)^{p}\, p^{-p}\, (1-p)^{-(1-p)}.
\]
\end{lemma}

\begin{lemma}[Exponential form of the $q$-ary entropy]
\lean{ErrorCorrectingCodes.Codeword.q_pow_qary_entropy_simp'}
\label{ErrorCorrectingCodes.Codeword.q_pow_qary_entropy_simp'}
\leanok
\uses{ErrorCorrectingCodes.Codeword.q_pow_qary_entropy_simp}
\difficulty{1}
Let $q \ge 2$ be an integer and let $p$ be a real number with $0 < p < 1$. Then the
base-$q$ power of the $q$-ary entropy $H_q(p)$ admits the closed form
\[
q^{\,H_q(p)} \;=\; (q-1)^{p}\, p^{-p}\,(1-p)^{-(1-p)}.
\]
\end{lemma}

%%%%%%%%%%%%%%%%%%%%%%%%%%%%%%%%%%%%%%%%%%%%%%%%%%%%%%%%%%%%%%%%%%%%%%%%%%%%%%%
\section{Analytic helpers}

\begin{lemma}[Square root minus square root of floor is at most one]
\lean{ErrorCorrectingCodes.Codeword.sqrt_sub_sqrt_floor_le_one}
\label{ErrorCorrectingCodes.Codeword.sqrt_sub_sqrt_floor_le_one}
\leanok
\difficulty{5}
For every real number $x \ge 0$,
\[
\sqrt{x} - \sqrt{\lfloor x \rfloor} \le 1,
\]
where $\lfloor x \rfloor$ denotes the greatest integer not exceeding $x$.
\end{lemma}

%%%%%%%%%%%%%%%%%%%%%%%%%%%%%%%%%%%%%%%%%%%%%%%%%%%%%%%%%%%%%%%%%%%%%%%%%%%%%%%
\section{Stirling-based binomial lower bound}

\begin{lemma}[Asymptotic lower bound on a weighted binomial coefficient]
\lean{ErrorCorrectingCodes.Codeword.binomial_coef_asymptotic_lower_bound'}
\label{ErrorCorrectingCodes.Codeword.binomial_coef_asymptotic_lower_bound'}
\leanok
\proofsource{ada-coding-book}{pdf-7251cb8e4bc5-p072-b001}
\uses{ErrorCorrectingCodes.Codeword.qaryEntropy}
\difficulty{7}
Fix a real number $p$ with $0 < p < 1$ and an integer $q \ge 2$, and let $H_q$ denote
the $q$-ary entropy function. Then there is a function $\varepsilon \colon \bbn \to
\bbr$ with $\varepsilon(n) = o(n)$ such that, for all sufficiently large $n$,
\[
\binom{n}{\lfloor np \rfloor}\,(q-1)^{pn} \;\ge\; q^{\,H_q(p)\,n \,-\, \varepsilon(n)}.
\]
\end{lemma}

%%%%%%%%%%%%%%%%%%%%%%%%%%%%%%%%%%%%%%%%%%%%%%%%%%%%%%%%%%%%%%%%%%%%%%%%%%%%%%%
\section{Positivity of $q$-ary entropy}

\begin{theorem}[Positivity of the $q$-ary entropy]
\lean{ErrorCorrectingCodes.Codeword.qary_entropy_pos}
\label{ErrorCorrectingCodes.Codeword.qary_entropy_pos}
\leanok
\uses{ErrorCorrectingCodes.Codeword.qaryEntropy}
\difficulty{5}
Let $\alpha$ be a finite alphabet with $q = \abs{\alpha}$ symbols, and let $p$ be a real
number with $0 < p \le 1 - 1/q$. Then the $q$-ary entropy of $p$ is strictly positive:
\[
H_q(p) \;=\; p\log_q(q-1) - p\log_q p - (1-p)\log_q(1-p) \;>\; 0.
\]
\end{theorem}

%%%%%%%%%%%%%%%%%%%%%%%%%%%%%%%%%%%%%%%%%%%%%%%%%%%%%%%%%%%%%%%%%%%%%%%%%%%%%%%
%%% added by the blueprint pipeline
\section{Additional declarations}

\begin{lemma}[Rearrangement of the $q$-ary entropy expression]
\lean{ErrorCorrectingCodes.Codeword.qary_entropy_logb_expand}
\label{ErrorCorrectingCodes.Codeword.qary_entropy_logb_expand}
\leanok
\difficulty{2}
For every natural number $q$ and every real number $p$,
\[
p\,\log_q(q-1) - p\,\log_q p - (1-p)\,\log_q(1-p)
=
\log_q(q-1)\cdot p + \log_q p\cdot(-p) + \log_q(1-p)\cdot\bigl(-(1-p)\bigr),
\]
where $\log_q$ denotes the base-$q$ logarithm. That is, the three terms of the $q$-ary
entropy expression may be rewritten by commuting each product and absorbing the
subtractions into the sign of the multiplier.
\end{lemma}

\begin{lemma}[AM–GM inequality for two nonnegative reals]
\lean{ErrorCorrectingCodes.Codeword.am_gm_sqrt_le_half_sum}
\label{ErrorCorrectingCodes.Codeword.am_gm_sqrt_le_half_sum}
\leanok
\difficulty{3}
For any nonnegative real numbers $a, b \ge 0$, the geometric mean is at most the
arithmetic mean:
\[
\sqrt{ab} \;\le\; \frac{a+b}{2}.
\]
\end{lemma}

\begin{lemma}[A Stirling square-root identity]
\lean{ErrorCorrectingCodes.Codeword.sqrt_two_pi_mul_sqrt_pi_half}
\label{ErrorCorrectingCodes.Codeword.sqrt_two_pi_mul_sqrt_pi_half}
\leanok
\difficulty{3}
For every real number $n \ge 0$, the three square roots $\sqrt{2\pi n}$, $\sqrt{\pi/2}$,
and $\sqrt{n}$ multiply to $\pi n$; that is,
\[
\sqrt{2\pi n}\,\cdot\,\sqrt{\pi/2}\,\cdot\,\sqrt{n} \;=\; \pi n.
\]
\end{lemma}

\begin{lemma}[Positivity of $\lfloor np\rfloor$ for large $n$]
\lean{ErrorCorrectingCodes.Codeword.floor_np_pos}
\label{ErrorCorrectingCodes.Codeword.floor_np_pos}
\leanok
\difficulty{4}
Let $p$ be a real number with $0 < p$ and $0 < 1 - p$, and set $N_2 = \lceil
2/(p(1-p))\rceil + 1$. Then for every natural number $n$ with $N_2 \le n$, the floor
$\lfloor np\rfloor$ is strictly positive, i.e. $\lfloor np\rfloor \ge 1$.
\end{lemma}

\begin{lemma}[Stirling-type lower bound for a factorial ratio]
\lean{ErrorCorrectingCodes.Codeword.stirling_comb_bound}
\proofsource{ada-coding-book}{pdf-7251cb8e4bc5-p072-b001}
\label{ErrorCorrectingCodes.Codeword.stirling_comb_bound}
\leanok
\proofstep
  {ErrorCorrectingCodes.Codeword}
  {context}
  {ada-coding-book}
  {pdf-7251cb8e4bc5-p016-b001,
    pdf-7251cb8e4bc5-p015-b006}
\proofstep
  {ErrorCorrectingCodes.Codeword.add}
  {context}
  {ada-coding-book}
  {pdf-7251cb8e4bc5-p042-b003,
    pdf-7251cb8e4bc5-p028-b004}
\proofstep
  {ErrorCorrectingCodes.Codeword.sub}
  {context}
  {ada-coding-book}
  {pdf-7251cb8e4bc5-p047-b004}
\proofstep
  {ErrorCorrectingCodes.Codeword.zero}
  {context}
  {ada-coding-book}
  {pdf-7251cb8e4bc5-p042-b005,
    pdf-7251cb8e4bc5-p029-b002}
\proofstep
  {ErrorCorrectingCodes.Codeword.hamming_distance}
  {direct}
  {ada-coding-book}
  {pdf-7251cb8e4bc5-p018-b002}
\proofstep
  {ErrorCorrectingCodes.Codeword.hamming_ball}
  {direct}
  {ada-coding-book}
  {pdf-7251cb8e4bc5-p029-b006}
\proofstep
  {ErrorCorrectingCodes.Codeword.am_gm_sqrt_le_half_sum}
  {direct}
  {ada-coding-book}
  {pdf-7251cb8e4bc5-p166-b003}
\proofstep
  {ErrorCorrectingCodes.Codeword.sqrt_two_pi_mul_sqrt_pi_half}
  {context}
  {ada-coding-book}
  {pdf-7251cb8e4bc5-p072-b001,
    pdf-7251cb8e4bc5-p201-b005}
\uses{ErrorCorrectingCodes.Codeword.am_gm_sqrt_le_half_sum, ErrorCorrectingCodes.Codeword.sqrt_two_pi_mul_sqrt_pi_half}
\difficulty{5}
Let $a,b,n\in\bbn$ with $a>0$, $b>0$, and $a+b=n$, and let $c>0$ be a real number.
Suppose the factorial product $a!\,b!$ admits the Stirling-type upper bound
\[
a!\,b! \;\le\;
c\cdot\Bigl(\sqrt{2}\,\sqrt{a}\,\sqrt{\pi}\,(a/e)^{a}\Bigr)\Bigl(\sqrt{2}\,\sqrt{b}\,\sqrt{\pi}\,(b/e)^{b}\Bigr).
\]
Then the ratio $n!/(a!\,b!)$ is bounded below by
\[
\frac{n^{n}}{a^{a}\,b^{b}}\cdot\frac{1}{c\,\sqrt{\pi/2}\,\sqrt{n}}
\;\le\;
\frac{n!}{a!\,b!}.
\]
\end{lemma}
